
\documentclass[addpoints]{exam}
\usepackage[version=4]{mhchem}
\usepackage{siunitx}
\usepackage{color}
\usepackage{amsmath}
\usepackage{booktabs}

% \printanswers

\newcommand{\event}{Final Exam}
\newcommand{\tournament}{}  % Replace with class name, date, or course code
\def\type{0}

\setlength\fillinlinelength{\textwidth}
\CorrectChoiceEmphasis{\color{red}\bfseries}
\SolutionEmphasis{\color{red}\bfseries}

\setlength\answerclearance{2pt}
\pagestyle{headandfoot}
\runningheadrule
\runningheader{\event}{\tournament}{\ifnum\type=1 Class Set Test \else Full Name:\hspace{1cm} \fi}
\runningfootrule
\runningfooter{}{Page \thepage\ of \numpages}{}

\begin{document}
\begin{center}
    {\huge Grade 12 Chemistry \\ \event
    \ifprintanswers
     \space Key
    \else
     \space \ifnum\type<2 Test \fi \ifnum\type=1 \textbf{(CLASS SET)} \fi \ifnum\type=2 Answer Sheet \fi
    \fi \\ \vspace{3mm}\tournament\vspace{1mm}}
\end{center}

\vspace{.25\textheight}

\begin{center}
    First Name: \ifprintanswers
    \underline{\hspace{3.5cm}}\textbf{KEY}\underline{\hspace{5.5cm}}\\\vspace{5mm}
    \else
        \underline{\hspace{10cm}}\\ \vspace{5mm}
    \fi
    Last Name: \underline{\hspace{10cm}}\\ \vspace{5mm}
\end{center}

\noindent\textbf{\underline{Directions:}}
\begin{itemize}
    \ifnum\type=1 \item \textbf{This test is a class set.} Please do not write any answers on this paper. \fi
    \item This exam covers all five units of Grade 12 Chemistry (SCH4U).
    \item Show all work for full marks on short answer questions.
    \item $K_w = 1.0\times10^{-14}$ at \SI{25}{\celsius};
          $F = \SI{9.65e4}{C/mol}$;
          $R = \SI{8.314}{J\,mol^{-1}K^{-1}}$.
    \item The exam is designed to be completed in 120 minutes.
\end{itemize}
\begin{center}
\textit{For grading use only}\\
\multirowgradetable{2}[pages]
\end{center}

\newpage
\begin{questions}

\section*{Part A --- Multiple Choice (20 marks)}
\vspace{3mm}
\noindent\textit{Questions 1--4: Structure and Properties of Matter}
\vspace{3mm}

\question[1] Which of the following is the correct ground-state electron configuration of
\ce{Cu} ($Z = 29$)?
\begin{choices}
    \choice $[\text{Ar}]3d^9 4s^2$
    \correctchoice $[\text{Ar}]3d^{10} 4s^1$
    \choice $[\text{Ar}]3d^8 4s^2$
    \choice $[\text{Ar}]3d^{10} 4s^2$
\end{choices}

\question[1] Which of the following molecules has a \textit{bent} molecular geometry?
\begin{choices}
    \choice \ce{CO2}
    \choice \ce{BeCl2}
    \correctchoice \ce{SO2}
    \choice \ce{CS2}
\end{choices}

\question[1] Which substance has the highest boiling point?
\begin{choices}
    \choice \ce{CH3OCH3} (dimethyl ether, $M = \SI{46}{g/mol}$)
    \correctchoice \ce{CH3CH2OH} (ethanol, $M = \SI{46}{g/mol}$)
    \choice \ce{C3H8} (propane, $M = \SI{44}{g/mol}$)
    \choice \ce{N2} ($M = \SI{28}{g/mol}$)
\end{choices}

\question[1] Which of the following best describes a metallic bond?
\begin{choices}
    \choice Electrostatic attraction between oppositely charged ions.
    \choice Covalent sharing of electrons between two non-metal atoms.
    \correctchoice Attraction between positive metal ions and a sea of delocalized electrons.
    \choice Hydrogen bonding between metal atoms.
\end{choices}

\vspace{3mm}
\noindent\textit{Questions 5--8: Organic Chemistry}
\vspace{3mm}

\question[1] What is the IUPAC name of the compound with the formula
\ce{CH3CH2CH(OH)CH3}?
\begin{choices}
    \choice Butan-1-ol
    \correctchoice Butan-2-ol
    \choice 2-butanol (not IUPAC)
    \choice Propan-2-ol
\end{choices}

\question[1] Which reaction type converts a carboxylic acid and an alcohol into an ester
and water?
\begin{choices}
    \choice Addition
    \choice Elimination
    \choice Substitution
    \correctchoice Condensation (esterification)
\end{choices}

\question[1] Which of the following polymers is formed by condensation polymerization?
\begin{choices}
    \choice Polyethylene
    \choice Polystyrene
    \correctchoice Nylon-6,6
    \choice Poly(tetrafluoroethylene) (PTFE)
\end{choices}

\question[1] An unknown compound reacts with \ce{Na} metal to produce \ce{H2} gas and gives
a positive Tollens' test (silver mirror). The compound is most likely a(n):
\begin{choices}
    \choice Ketone
    \correctchoice Hydroxy aldehyde (aldehyde with an OH group)
    \choice Ester
    \choice Alkene
\end{choices}

\vspace{3mm}
\noindent\textit{Questions 9--12: Energy Changes and Rates of Reaction}
\vspace{3mm}

\question[1] The combustion of propane is highly exothermic. Which statement is true?
\begin{choices}
    \correctchoice The total bond energy of the products is greater than the total bond
    energy of the reactants.
    \choice The total bond energy of the reactants is greater than the total bond energy
    of the products.
    \choice $\Delta H_{comb} > 0$.
    \choice The activation energy equals $\Delta H$.
\end{choices}

\question[1] The rate law for a reaction is rate $= k[\text{A}]^2$. If $[\text{A}]$ is
tripled, the rate will:
\begin{choices}
    \choice Triple.
    \choice Double.
    \correctchoice Increase by a factor of 9.
    \choice Increase by a factor of 6.
\end{choices}

\question[1] Which of the following changes \textit{does not} alter the value of the rate
constant $k$?
\begin{choices}
    \choice Increasing the temperature.
    \choice Adding a catalyst.
    \correctchoice Increasing the concentration of a reactant.
    \choice Changing the solvent.
\end{choices}

\question[1] For the following reactions, which can be used with Hess's Law to find
$\Delta H$ for \ce{C(s) + O2(g) -> CO2(g)}?
\begin{choices}
    \choice Only reactions with positive $\Delta H$.
    \choice Only reactions with the same number of moles of gas.
    \correctchoice Any combination of reactions that sums to the target equation.
    \choice Only reactions at the same temperature.
\end{choices}

\vspace{3mm}
\noindent\textit{Questions 13--16: Chemical Systems and Equilibrium}
\vspace{3mm}

\question[1] For the equilibrium \ce{2NO2(g) <=> N2O4(g)}, if $Q > K_{eq}$, the system will:
\begin{choices}
    \choice Produce more \ce{NO2}.
    \correctchoice Convert \ce{N2O4} to \ce{NO2} until equilibrium is re-established.
    \choice Remain unchanged.
    \choice Increase $K_{eq}$.
\end{choices}

\question[1] For the reaction \ce{PCl5(g) <=> PCl3(g) + Cl2(g)}, $\Delta H > 0$.
Increasing temperature will:
\begin{choices}
    \correctchoice Shift equilibrium right and increase $K_{eq}$.
    \choice Shift equilibrium left and decrease $K_{eq}$.
    \choice Shift equilibrium right and decrease $K_{eq}$.
    \choice Have no effect because $K_{eq}$ does not depend on temperature.
\end{choices}

\question[1] What is the pH of a \SI{0.050}{mol/L} solution of \ce{Ba(OH)2}, a strong base
that fully dissociates to give 2 \ce{OH-} per formula unit?
\begin{choices}
    \choice 12.3
    \correctchoice 13.0
    \choice 11.7
    \choice 1.0
\end{choices}

\question[1] Which of the following will \textit{increase} the molar solubility of \ce{BaSO4}
($K_{sp} = 1.1\times10^{-10}$)?
\begin{choices}
    \choice Adding more \ce{BaSO4(s)} to the solution.
    \choice Adding \ce{Na2SO4(aq)}.
    \choice Decreasing the temperature (assuming $K_{sp}$ decreases).
    \correctchoice Adding \ce{HCl(aq)} (which reacts with \ce{SO4^{2-}}).
\end{choices}

\vspace{3mm}
\noindent\textit{Questions 17--20: Electrochemistry}
\vspace{3mm}

\question[1] In an electrolytic cell, which statement is correct?
\begin{choices}
    \choice Oxidation occurs at the cathode.
    \choice The reaction is spontaneous.
    \correctchoice An external power source drives a non-spontaneous reaction.
    \choice Electrons flow from cathode to anode through the external circuit.
\end{choices}

\question[1] Given $E^\circ(\ce{Al^{3+}/Al}) = -\SI{1.66}{V}$ and
$E^\circ(\ce{Sn^{2+}/Sn}) = -\SI{0.14}{V}$, the standard cell potential for an
Al--Sn galvanic cell is:
\begin{choices}
    \choice $-\SI{1.52}{V}$
    \correctchoice $+\SI{1.52}{V}$
    \choice $+\SI{1.80}{V}$
    \choice $-\SI{1.80}{V}$
\end{choices}

\question[1] How many moles of electrons are required to deposit \SI{1.00}{mol} of
\ce{Al} from \ce{Al^{3+}}?
\begin{choices}
    \choice 1
    \choice 2
    \correctchoice 3
    \choice 6
\end{choices}

\question[1] In the reaction \ce{Sn^{2+}(aq) + 2Fe^{3+}(aq) -> Sn^{4+}(aq) + 2Fe^{2+}(aq)},
which species is oxidized?
\begin{choices}
    \choice \ce{Fe^{3+}}
    \choice \ce{Fe^{2+}}
    \correctchoice \ce{Sn^{2+}}
    \choice \ce{Sn^{4+}}
\end{choices}


\section*{Part B --- Short Answer (40 marks)}
\vspace{3mm}
\noindent\textit{Answer all five questions.}
\vspace{3mm}

\question \textbf{(Structure and Properties of Matter)}
\begin{parts}
    \part[3] Phosphorus pentachloride (\ce{PCl5}) reacts with water to form phosphoric acid and
    hydrogen chloride gas:
    \[
    \ce{PCl5(g) + 4H2O(l) -> H3PO4(aq) + 5HCl(g)}
    \]
    For the \ce{PCl5} molecule, determine: (i) the hybridization of P, (ii) the molecular
    geometry, and (iii) whether the molecule is polar or nonpolar.
    \part[3] Compare the following three substances in terms of their dominant intermolecular
    forces and predict their relative boiling points, ranking from lowest to highest:
    \ce{F2}, \ce{HF}, \ce{NaF}.
    \part[2] The ionic compound \ce{MgO} has a much higher melting point (\SI{2852}{\celsius})
    than \ce{NaCl} (\SI{801}{\celsius}). Explain this difference using Coulomb's law and
    lattice energy.
\end{parts}
\begin{solution}[10cm]
\textbf{(a)} \ce{PCl5}: (i) \textbf{$sp^3d$} hybridization (5 bonding pairs, expanded octet).
(ii) Molecular geometry: \textbf{trigonal bipyramidal}. (iii) \textbf{Nonpolar} --- the five
P--Cl bond dipoles are arranged symmetrically and cancel completely.

\textbf{(b)}
\begin{itemize}
    \item \ce{F2}: \textbf{London dispersion} only (nonpolar, diatomic). Very low boiling point
    (\SI{-188}{\celsius}).
    \item \ce{HF}: \textbf{London dispersion, dipole--dipole, hydrogen bonding} (strong H-bond).
    Higher boiling point (\SI{19}{\celsius}).
    \item \ce{NaF}: \textbf{Ionic bonding} (not an intermolecular force but an ionic lattice).
    Highest melting/boiling point (\SI{993}{\celsius}).
\end{itemize}
Ranking: \ce{F2} $<$ \ce{HF} $<$ \ce{NaF}.

\textbf{(c)} Lattice energy $\propto \dfrac{q_+\,q_-}{r}$. In \ce{MgO}: $q_+ = +2$,
$q_- = -2$, $r(\ce{MgO}) < r(\ce{NaCl})$ (smaller ions). In \ce{NaCl}: $q_+ = +1$,
$q_- = -1$. The product of charges in \ce{MgO} is $2\times2 = 4$ vs.\ $1\times1 = 1$ for
\ce{NaCl}. The much larger charges (and smaller ionic radius) give \ce{MgO} a far greater
lattice energy, requiring much more energy to melt.
\end{solution}

\question \textbf{(Organic Chemistry)}

An organic compound A has the molecular formula \ce{C4H8O2}.
\begin{itemize}
    \item It has a pleasant fruity odour.
    \item It does not react with \ce{Na} metal.
    \item It can be hydrolysed in aqueous \ce{NaOH} to produce an alcohol and a carboxylate
    salt.
\end{itemize}
\begin{parts}
    \part[2] Identify the functional group and class of compound A. Name compound A if it is
    ethyl ethanoate.
    \part[3] Write the equation for the hydrolysis of ethyl ethanoate (\ce{CH3COOC2H5}) in
    aqueous \ce{NaOH}. Identify the products and name the reaction type.
    \part[3] Starting from ethanol and ethanoic acid, describe in words the steps and
    conditions required to synthesize ethyl ethanoate. Write the balanced equation and explain
    the role of the \ce{H2SO4} catalyst.
\end{parts}
\begin{solution}[10cm]
\textbf{(a)} Functional group: \textbf{ester} (\ce{-COO-}). Does not react with Na (no OH),
pleasant fruity odour and hydrolysable confirm it is an ester. Compound A is
\textbf{ethyl ethanoate} (\ce{CH3COOC2H5}).

\textbf{(b)} Saponification (base hydrolysis):
\[
\ce{CH3COOC2H5(l) + NaOH(aq) -> CH3COO-Na+(aq) + C2H5OH(aq)}
\]
Products: \textbf{sodium ethanoate} (carboxylate salt) and \textbf{ethanol} (alcohol).
Reaction type: \textbf{saponification (hydrolysis / condensation reverse)}.

\textbf{(c)} Esterification conditions: mix ethanol and ethanoic acid with a few drops of
concentrated \ce{H2SO4} catalyst; heat gently under reflux.
\[
\ce{CH3COOH(l) + C2H5OH(l) <=>[H2SO4][\Delta] CH3COOC2H5(l) + H2O(l)}
\]
The \ce{H2SO4} catalyst protonates the carbonyl oxygen of the carboxylic acid, activating it
toward nucleophilic attack by the alcohol. It lowers the activation energy and speeds up the
reaction without shifting the equilibrium position (does not change yield). The reaction is
reversible, so excess of one reactant or removal of water improves yield.
\end{solution}

\question \textbf{(Energy Changes and Rates of Reaction)}
\begin{parts}
    \part[4] \SI{50.0}{mL} of \SI{1.00}{mol/L} \ce{HCl} is mixed with \SI{50.0}{mL} of
    \SI{1.00}{mol/L} \ce{NaOH} in a coffee-cup calorimeter. The temperature rises from
    \SI{21.5}{\celsius} to \SI{28.3}{\celsius}. Assume the solution density is
    \SI{1.00}{g/mL} and $c = \SI{4.18}{J\,g^{-1}\,\celsius^{-1}}$.
    Calculate $\Delta H_{neutralization}$ in kJ per mole of water formed.
    \part[4] Using Hess's Law, calculate $\Delta H$ for the reaction:
    \[
    \ce{2C(s) + H2(g) -> C2H2(g)} \quad \Delta H = \, ?
    \]
    given:
    \begin{align*}
    \text{(1)}\; & \ce{C2H2(g) + 5/2 O2(g) -> 2CO2(g) + H2O(l)} \; \Delta H_1 = -\SI{1300}{kJ}\\
    \text{(2)}\; & \ce{C(s) + O2(g) -> CO2(g)} \; \Delta H_2 = -\SI{394}{kJ}\\
    \text{(3)}\; & \ce{H2(g) + 1/2 O2(g) -> H2O(l)} \; \Delta H_3 = -\SI{286}{kJ}
    \end{align*}
\end{parts}
\begin{solution}[10cm]
\textbf{(a)} Total volume $= \SI{100.0}{mL}$; total mass $= \SI{100.0}{g}$.
\[
q_{soln} = mc\Delta T = (100.0)(4.18)(28.3 - 21.5) = (100.0)(4.18)(6.8) = \SI{2842}{J}
\]
$q_{rxn} = -q_{soln} = \SI{-2842}{J} = \SI{-2.842}{kJ}$ (exothermic).

Moles of \ce{HCl} = moles of \ce{NaOH} = $(0.0500\ \text{L})(1.00\ \text{mol/L}) =
\SI{0.0500}{mol}$. One mole of \ce{H2O} is formed per mole of acid.
\[
\Delta H_{neutralization} = \frac{-2.842}{0.0500} = \mathbf{-56.8\ kJ/mol}
\]

\textbf{(b)} To obtain \ce{2C(s) + H2(g) -> C2H2(g)}, manipulate the given equations:
\begin{itemize}
    \item Reverse (1): \ce{2CO2(g) + H2O(l) -> C2H2(g) + 5/2 O2(g)}; $\Delta H = +1300\ \text{kJ}$
    \item Keep (2) $\times$ 2: $2 \times (-394) = -788\ \text{kJ}$
    \item Keep (3): $\Delta H_3 = -286\ \text{kJ}$
\end{itemize}
Adding and cancelling (\ce{2CO2}, \ce{H2O}, \ce{O2}):
\[
\Delta H = +1300 + (-788) + (-286) = \mathbf{+226\ kJ}
\]
(Endothermic, consistent with the formation of an unstable triple-bonded molecule.)
\end{solution}

\question \textbf{(Chemical Systems and Equilibrium)}

Nitrous acid (\ce{HNO2}) is a weak acid with $K_a = 4.5\times10^{-4}$ at \SI{25}{\celsius}.
\begin{parts}
    \part[4] Calculate the pH and percent dissociation of a \SI{0.200}{mol/L} \ce{HNO2}
    solution. Construct a complete ICE table.
    \part[4] A buffer solution is prepared by mixing \SI{0.200}{mol/L} \ce{HNO2} and
    \SI{0.150}{mol/L} \ce{NaNO2}.
    \begin{itemize}
        \item[(i)] Calculate the pH of this buffer using the Henderson--Hasselbalch equation:
        $\text{pH} = pK_a + \log\dfrac{[\text{A}^-]}{[\text{HA}]}$
        \item[(ii)] Explain qualitatively what happens to the pH when a small amount of
        \ce{HCl} is added to this buffer. Write the equation for the reaction that occurs.
    \end{itemize}
\end{parts}
\begin{solution}[10cm]
\textbf{(a)} ICE table for \ce{HNO2(aq) <=> H+(aq) + NO2-(aq)}:

\begin{center}
\begin{tabular}{lccc}
\toprule
 & \ce{HNO2} & \ce{H+} & \ce{NO2-} \\
\midrule
I & 0.200 & 0 & 0 \\
C & $-x$ & $+x$ & $+x$ \\
E & $0.200-x$ & $x$ & $x$ \\
\bottomrule
\end{tabular}
\end{center}

\[
K_a = \frac{x^2}{0.200 - x} = 4.5\times10^{-4}
\]
Assume $x \ll 0.200$:
\[
x^2 = (4.5\times10^{-4})(0.200) = 9.0\times10^{-5} \implies x = 9.49\times10^{-3}\ \text{mol/L}
\]
\[
\text{pH} = -\log(9.49\times10^{-3}) = \mathbf{2.02}
\]
\[
\text{\% dissociation} = \frac{9.49\times10^{-3}}{0.200}\times100\% = \mathbf{4.7\%}
\]
(Check: $4.7\% < 5\%$, assumption valid.)

\textbf{(b)(i)}
\[
pK_a = -\log(4.5\times10^{-4}) = 3.35
\]
\[
\text{pH} = 3.35 + \log\frac{0.150}{0.200} = 3.35 + \log(0.750) = 3.35 + (-0.125) = \mathbf{3.22}
\]

\textbf{(b)(ii)} When \ce{HCl} is added, \ce{H+} ions are introduced. The conjugate base
\ce{NO2-} in the buffer reacts with the added \ce{H+}:
\[
\ce{NO2-(aq) + H+(aq) -> HNO2(aq)}
\]
This consumes most of the added acid and converts it to the weak acid, so the pH decreases
only slightly. Without the buffer, the same amount of \ce{HCl} would cause a large drop in pH.
\end{solution}

\question \textbf{(Electrochemistry)}

Consider the following standard reduction potentials:
\[
E^\circ(\ce{Cu^{2+}/Cu}) = +\SI{0.34}{V}, \quad E^\circ(\ce{Fe^{2+}/Fe}) = -\SI{0.44}{V}
\]
\begin{parts}
    \part[2] Write the spontaneous overall cell reaction and calculate $E^\circ_{cell}$.
    \part[2] Calculate $\Delta G^\circ$ for the cell reaction.
    ($F = \SI{9.65e4}{C/mol}$)
    \part[4] Copper metal is to be electroplated from a \ce{CuSO4(aq)} solution onto an iron
    object.
    \begin{itemize}
        \item[(i)] In the electrolytic cell for this process, identify which electrode is the
        anode and which is the cathode.
        \item[(ii)] Write the half-reaction at the cathode.
        \item[(iii)] Calculate the mass of Cu deposited if a current of \SI{3.00}{A} flows
        for \SI{30.0}{min}. ($M_{\ce{Cu}} = \SI{63.55}{g/mol}$)
    \end{itemize}
\end{parts}
\begin{solution}[10cm]
\textbf{(a)} Lower reduction potential = anode (Fe is oxidized):
\begin{align*}
\text{Anode (ox):}\  & \ce{Fe(s) -> Fe^{2+}(aq) + 2e-} \\
\text{Cathode (red):}\  & \ce{Cu^{2+}(aq) + 2e- -> Cu(s)}
\end{align*}
Overall: $\ce{Fe(s) + Cu^{2+}(aq) -> Fe^{2+}(aq) + Cu(s)}$
\[
E^\circ_{cell} = E^\circ_{cathode} - E^\circ_{anode} = (+0.34) - (-0.44) = \mathbf{+0.78\ V}
\]

\textbf{(b)} $n = 2$ electrons transferred.
\[
\Delta G^\circ = -nFE^\circ_{cell} = -(2)(9.65\times10^4)(0.78) = \mathbf{-1.51\times10^5\ J
= -151\ kJ}
\]
Reaction is spontaneous ($\Delta G^\circ < 0$, $E^\circ_{cell} > 0$).

\textbf{(c)}
\begin{itemize}
    \item[(i)] In the electroplating cell: the \textbf{iron object is the cathode} (copper
    is deposited on it); the \textbf{copper electrode or inert anode} is the anode.
    \item[(ii)] Cathode: \ce{Cu^{2+}(aq) + 2e- -> Cu(s)}
    \item[(iii)] $Q = It = (3.00)(30.0 \times 60) = (3.00)(1800) = \SI{5400}{C}$
    \[
    n_{e^-} = \frac{5400}{9.65\times10^4} = 5.596\times10^{-2}\ \text{mol}
    \]
    \[
    n_{\ce{Cu}} = \frac{5.596\times10^{-2}}{2} = 2.798\times10^{-2}\ \text{mol}
    \]
    \[
    m_{\ce{Cu}} = (2.798\times10^{-2})(63.55) = \mathbf{1.78\ g}
    \]
\end{itemize}
\end{solution}

\end{questions}
\end{document}
